% Zumdahl Chapter 11 free-response set -- 2 long (10) + 3 short (4) = 32 pts.
\documentclass{apchem-exam}

\apcunitword{Chapter}
\apcunit{11}
\apcunittitle{Properties of Solutions}
\apcdoctitle{Free-Response Set}
\apcsubtitle{Zumdahl \S11.1--\S11.3 \textbullet\ 5 questions \textbullet\ 32 points \textbullet\ 50 minutes}

\begin{document}
\apctitleblock

\begin{apcnote}[Scope --- most of this chapter is excluded]
Only \S11.1--11.3 are examinable, as \textbf{CED 3.7} (solutions and
mixtures) and \textbf{3.10} (solubility). Read this before assigning:

The CED exclusion on topic 3.8 removes \textbf{colligative properties
entirely}, and with them \textbf{calculations of molality, percent by
mass and percent by volume}. That takes \S11.4--11.8 --- vapour-pressure
lowering, boiling-point elevation, freezing-point depression, osmotic
pressure, the van't Hoff factor and colloids --- off the exam completely.
None of it appears below.

Concentration on the AP exam means \textbf{molarity}. Every calculation
here uses it.
\end{apcnote}

\examsection{II}{Free Response}{5 questions \textbullet\ 32 points \textbullet\ 50 minutes}{\calculatorok}

% =====================================================================
\frq{long}{10}{Zumdahl \S11.2 \textbullet\ CED 3.10 \textbullet\ the energetics of dissolving}

Dissolving an ionic solid in water can be broken into steps: separating
the ions from the lattice, and hydrating them.

\begin{parts}
  \item State the sign of the enthalpy change for each step, with a
        reason. \answerlines[3]{\textbf{Separating the lattice:
        endothermic, positive.} Energy must be supplied to overcome the
        Coulombic attractions holding the ions together.
        \textbf{Hydration: exothermic, negative.} Energy is released as
        water molecules are attracted to the ions and surround them.}
  \item Explain how the relative sizes of those two terms determine
        whether dissolving is exothermic or endothermic overall.
        \answerlines[3]{The overall $\Delta H_{\text{soln}}$ is their sum.
        If hydration releases more than lattice separation costs, the
        total is negative and the solution warms. If lattice separation
        costs more, the total is positive and the solution cools. The two
        terms are individually large and often nearly cancel, which is why
        $\Delta H_{\text{soln}}$ is usually small.}
  \item Ammonium nitrate dissolves readily in water even though the
        solution becomes noticeably colder. Explain how a process that
        absorbs energy can still occur spontaneously.
        \workspace[3cm]
        \keyonly{\begin{apcanswer}Enthalpy is not the only factor.
        Dissolving scatters the ions, which were locked in an ordered
        crystal, throughout the solvent --- a large increase in the
        \textbf{disorder}, or entropy, of the system.

        That entropy increase is enough to drive the process even though
        it is endothermic. Spontaneity depends on both the energy change
        and the entropy change, not on enthalpy
        alone.\end{apcanswer}}
  \item Describe how water molecules arrange themselves around a dissolved
        \ce{Na+} ion and around a dissolved \ce{Cl-} ion.
        \answerlines[3]{Around \ce{Na+} the water molecules turn their
        \emph{oxygen} ends inward, because oxygen carries the partial
        negative charge and is attracted to the cation. Around \ce{Cl-}
        they turn their \emph{hydrogen} ends inward, since those carry the
        partial positive charge. In both cases the attraction is
        ion--dipole.}
\end{parts}

\begin{rubric}
  \rpoint{2}{(a) 1 pt lattice separation endothermic with a Coulombic
    reason; 1 pt hydration exothermic}
  \rpoint{3}{(b) 1 pt the total is the sum; 1 pt hydration dominant gives
    exothermic; 1 pt lattice dominant gives endothermic}
  \rpoint{3}{(c) 1 pt naming entropy / increased disorder; 1 pt the ions
    become dispersed from an ordered lattice; 1 pt spontaneity depends on
    more than enthalpy. ``It just happens'' earns 0}
  \rpoint{2}{(d) 1 pt correct orientation for each ion; 1 pt naming the
    ion--dipole interaction}
\end{rubric}

% =====================================================================
\frq{long}{10}{Zumdahl \S11.3 \textbullet\ CED 3.10 \textbullet\ factors affecting solubility}

\begin{parts}
  \item Explain, using ``like dissolves like'', why sodium chloride
        dissolves readily in water but not in hexane, \ce{C6H14}.
        \workspace[3cm]
        \keyonly{\begin{apcanswer}Water is \textbf{polar} and can surround
        ions with strong ion--dipole attractions. The energy released on
        hydration is large enough to repay the cost of pulling the ionic
        lattice apart, so \ce{NaCl} dissolves.

        Hexane is \textbf{non-polar}. Its only interactions are weak
        London dispersion forces, which cannot compensate for the very
        large lattice energy of \ce{NaCl}. There is no favourable
        interaction available to stabilize the separated ions, so the
        solid does not dissolve.\end{apcanswer}}
  \item The solubility of most solid solutes rises with temperature, but
        the solubility of a gas in water \emph{falls}. Explain the gas
        behaviour. \answerlines[3]{Dissolved gas molecules are held in
        solution only weakly. Raising the temperature increases their
        kinetic energy, so more of them have enough energy to escape the
        liquid, and the equilibrium shifts back toward the gas phase.
        Dissolving a gas is usually exothermic, so heating shifts it in
        the endothermic direction --- out of solution.}
  \item A sealed bottle of carbonated water fizzes when opened. Explain in
        terms of the pressure dependence of gas solubility.
        \answerlines[3]{The solubility of a gas is proportional to its
        partial pressure above the liquid. Sealed, the bottle holds
        \ce{CO2} at high pressure, so a large amount stays dissolved.
        Opening it drops the partial pressure of \ce{CO2} sharply, the
        solution is now supersaturated, and the excess gas comes out of
        solution as bubbles.}
  \item State whether pressure has a comparable effect on the solubility
        of a solid, and why. \answerlines[2]{No. Solids and liquids are
        nearly incompressible, so changing the pressure barely changes
        their volume or the equilibrium between dissolved and undissolved
        solute. The effect is negligible.}
\end{parts}

\begin{rubric}
  \rpoint{3}{(a) 1 pt water polar with ion--dipole attraction; 1 pt hexane
    non-polar with only dispersion; 1 pt the comparison with lattice
    energy}
  \rpoint{2}{(b) 1 pt higher kinetic energy lets gas molecules escape;
    1 pt equilibrium shifts out of solution}
  \rpoint{3}{(c) 1 pt solubility proportional to partial pressure; 1 pt
    pressure drops on opening; 1 pt excess gas leaves as bubbles}
  \rpoint{2}{(d) 1 pt negligible effect; 1 pt because solids are
    essentially incompressible}
\end{rubric}

% =====================================================================
\frq{short}{4}{Zumdahl \S11.1 \textbullet\ CED 3.7 \textbullet\ preparing a solution of known molarity}

\begin{parts}
  \item Calculate the mass of \ce{NaCl} needed to prepare
        \SI{250.0}{\milli\litre} of a \SI{0.150}{\Molar} solution.
        \workspace[2.8cm]
        \keyonly{\begin{apcanswer}$n = MV = (0.150)(0.2500) =
        \num{0.0375}$~mol

        $m = (0.0375)(58.44) = \mathbf{2.19}$~g\end{apcanswer}}
  \item Describe why the solid is dissolved in \emph{less} than
        \SI{250}{\milli\litre} of water and the solution then made up to
        the mark, rather than being added to
        \SI{250}{\milli\litre} of water. \answerlines[3]{Molarity is
        defined per litre of \emph{solution}, not per litre of solvent.
        The dissolved solid changes the total volume, so adding the solid
        to \SI{250}{\milli\litre} of water would give a final volume
        greater than \SI{250}{\milli\litre} and a concentration below
        \SI{0.150}{\Molar}.}
\end{parts}

\begin{rubric}
  \rpoint{2}{(a) 1 pt \num{0.0375}~mol; 1 pt \SI{2.19}{\gram}}
  \rpoint{2}{(b) 1 pt molarity is per litre of solution; 1 pt the solute
    contributes to the final volume}
\end{rubric}

% =====================================================================
\frq{short}{4}{Zumdahl \S11.3 \textbullet\ CED 3.10 \textbullet\ structure and solubility}

Ethanol (\ce{CH3CH2OH}) is miscible with water in all proportions.
Hexane is not.

\begin{parts}
  \item Explain the difference in terms of intermolecular forces.
        \answerlines[3]{Ethanol has an \ce{-OH} group, so it can form
        \textbf{hydrogen bonds} with water --- the new solute--solvent
        attractions are comparable in strength to the water--water
        attractions they replace. Hexane is a non-polar hydrocarbon
        limited to weak dispersion forces, which cannot compensate for
        breaking water's hydrogen-bond network.}
  \item Predict, with a reason, whether 1-octanol
        (\ce{CH3(CH2)7OH}) is as soluble in water as ethanol.
        \answerlines[3]{Much less soluble. It still has one \ce{-OH}
        group, but it carries a long non-polar hydrocarbon chain that
        dominates the molecule. The chain cannot interact favourably with
        water, so as the non-polar portion grows the solubility falls.}
\end{parts}

\begin{rubric}
  \rpoint{2}{(a) 1 pt ethanol hydrogen bonds with water via \ce{-OH};
    1 pt hexane offers only dispersion}
  \rpoint{2}{(b) 1 pt less soluble; 1 pt the long non-polar chain
    outweighs the single \ce{-OH}}
\end{rubric}

% =====================================================================
\frq{short}{4}{Zumdahl \S11.1 \textbullet\ CED 3.7 \textbullet\ concentration of ions}

A \SI{0.200}{\Molar} solution of \ce{Al2(SO4)3} is prepared and is fully
dissociated.

\begin{parts}
  \item Calculate the concentration of aluminium ions and of sulfate ions.
        \answerlines[3]{Each formula unit gives two \ce{Al^3+} and three
        \ce{SO4^2-}.

        $[\ce{Al^3+}] = 2(0.200) = \mathbf{0.400}$~M \quad
        $[\ce{SO4^2-}] = 3(0.200) = \mathbf{0.600}$~M}
  \item Calculate the total concentration of ions in the solution.
        \answerlines[2]{$0.400 + 0.600 = \mathbf{1.00}$~M, i.e.\ five
        times the concentration of the formula unit.}
\end{parts}

\begin{rubric}
  \rpoint{2}{(a) 1 pt each ion, using the correct subscript as a
    multiplier}
  \rpoint{2}{(b) 1 pt \SI{1.00}{\Molar}; 1 pt recognizing the factor of
    five from the formula. Answering \SI{0.200}{\Molar} earns 0}
\end{rubric}

\examtotal

\end{document}
